\item We remove the existential free-lift condition by an explicit
construction. Let $d>1$ be the order of $M$ and choose a primitive
$u\in\ker(M-I)\cap\mathbb Z^6$. Such a vector is obtained from a
nonzero rational kernel vector by clearing denominators and dividing
by the greatest common divisor of its integer coordinates; both
operations retain the equation $(M-I)u=0$. Put
$g(x)=Mx+u/d^2$ and $S_d=\sum_{a=0}^{d-1}M^a$.
Then $g^j(x)=M^jx+ju/d^2$, $S_dM^j=S_d$ and $S_du=du$.
A fixed point of $g^j$ gives
$(M^j-I)x+ju/d^2=\lambda\in\mathbb Z^6$.
Applying $S_d$ gives $ju/d=S_d\lambda\in\mathbb Z^6$.
A Bezout row pairing to one with $u$ therefore gives $d\mid j$.
Writing $j=dr$, the original equation becomes $ru/d=\lambda$,
so the same row gives $d\mid r$. Thus only multiples of $d^2$
can have fixed points. Since $g^{d^2}$ is translation by the
integer vector $u$, its torus action is the identity. This proves
both its actual order $d^2$ and freeness of every nonidentity power.
For $M=I$, translation by $u/2$ for primitive $u$ has actual order
two and is free. Conversely, when $I-M$ is invertible over
$\mathbb R$, every affine map $Mx+a$ fixes the torus class of
$(I-M)^{-1}a$. The free-lift test is consequently exactly
$\ker(M-I)\ne0$.

\item For an admissible full tuple $T=(M_1,\ldots,M_k,M_0)$ let
$H(T)$ be its maximum absolute entry. There is at most one identity
finite slot in a three-point tuple: two such identities would make
$M_0$ finite order, whereas $M_0=I+N$, $N^2=0$ and
$\operatorname{rank}N=3$ give $(I+N)^d=I+dN\ne I$ for every
positive integer $d$. At a fixed identity slot, delete that matrix
and reinsert $I$ at the same slot after classifying the remaining
ordered tuple. These two maps are inverse, and both commute with
$T\mapsto (QM_1Q^{-1},\ldots,QM_0Q^{-1})$ for the same
$Q\in\mathrm{GL}(6,\mathbb Z)$. They therefore retain the actual
returned simultaneous conjugator. The remaining invertible integral
matrices have height at least one, so inserting $I$ leaves height
unchanged. Every entry in the inserted block is constant; the first
unequal row-major entry of two tuples is therefore the same before
and after insertion. This proves preservation of the specified
lexicographic comparison as well. The actual affine order and
numerator at that slot are retained as gluing data and are transported
by this same $Q$. In the two-involution sector use the canonical
representative of \ref{maps:src:F.2.38}, with its own specified
convention. In all remaining sectors, order full integer tuples
first by $H$, then lexicographically by their row-major entries in
the displayed marked order, and take the first admissible member
of the simultaneous conjugacy class.
